\chapter{A full-parent remote theorem for strong HI sources}
\label{component:pcee1372511ed}
\noindent\textbf{Component }\path{pcee1372511ed}\par
{\small\noindent Source: \path{R13/baseline/baseline/advances/hi/remote_all_parent.tex}\par}
{\small\noindent SHA256: \path{cee1372511ed2fb356229e5f15bf95a827ce2829e9f8a51a88fe07db4fe72c28}\par}
{\small\noindent Inventory use: Every remote original parent state/angle chart and literal critical axis; only nonzero reciprocal-u and vInf corner remains\par}
\medskip
\section{A full-parent remote theorem for strong HI sources}
\label{pcee1372511ed:sec:hi-remote}
\begin{theorem}[All tests outside the HI critical corner]
Let $K_c,K_t$ be independent real closed extensions of $\mathbb R$.
Consider the original C2 parent with $d=1$, $\lambda=\beta=0$,
$n,A\geq0$, and $m,C,D,\alpha$ positive unbounded source numbers.
Then at every projective testing pair except possibly
$1/u$ and $v$ both nonzero infinitesimal, and every real or complex testing
state, the
original parent is a sum of squares in $K_c\otimes_{\mathbb R}K_t$.
No comparison of source and testing rates is made.
\end{theorem}

\begin{proof}
For any positive real $\mu$, retain the four nonnegative event
packets with coefficients $m-\mu,C-\mu,D-\mu,\alpha-\mu$, and
retain the $n$ and $A$ packets in their entirety. It suffices to
prove the capped parent obtained by setting
$m=C=D=\alpha=\mu$ and $n=A=0$ is tensor SOS at the given test.
The real constant $\mu$ may be chosen after the real testing
residue is fixed, since each of the original four source
coefficients is larger than every real number.

Write $D_u=1+u^2$, $D_v=1+v^2$.
First suppose $v$ is infinitesimal, including zero. Use the roots
\[
 S=vW,\qquad E=Y-uW,\qquad X,\quad Z.
\]
All four capped event amplitudes are bounded linear forms in
these roots:
\[
 M=D_uS-X-uZ+uvE,\quad uW-Y=-E,\quad
 J=-E+uvX-vZ,\quad vY-Z=uS+vE-Z.
\]
The remaining core depends only on the same four roots, so the
exact capped parent has a bounded four-root Gram.
Let $a=\bar u$ and $d_a=1+a^2$. Its residue is
\[
 (8\mu+2)E^2+
 4\mu(d_aS-X-aZ)^2+(4\mu+2)d_a(aS-Z)^2
       +2(X+aZ)^2-2d_aZ^2.
\]
Put $T=X+aZ$ and $R=aS-Z$; this is an invertible real change of
coordinates. The last three-variable form is
\[
 4\mu(d_aS-T)^2+4\mu d_aR^2+2T^2
       -2d_a a^2S^2+4d_a aSR.
\]
The $T$ and $R$ pivots are $4\mu+2$ and $4\mu d_a$.
The remaining Schur pivot is exactly
\[
 \frac{d_a[4\mu^2-a^2(4\mu+1)]}{\mu(2\mu+1)}.
\]
Taking $\mu=1+a^2$ makes its numerator positive, because
$4\mu^2-a^2(4\mu+1)=3a^2+4$.
Thus this bounded Gram has strictly positive real residue and
is tensor SOS by the inherited strict-residue Gram lemma.

Next suppose $v$ is a bounded unit. At the real residues
$(a,b)=(\bar u,\bar v)$ the four event forms have common kernel
\[
 (W,X,Y,Z)=W(1,b,a,ab).
\]
The core restricted to this line is
$2(1+a^2)(1+b^2)b^2W^2$, which is strictly positive away from zero.
Consequently a sufficiently large positive real $\mu$ makes the
whole capped residue form positive definite. For completeness,
this elementary assertion has a direct quantitative proof.
Let $E_0\succeq0$ be the event Gram, split the state space into
$\ker E_0$ and its orthogonal complement, and let $\sigma>0$ be
its smallest eigenvalue on that complement. Let $\kappa>0$ be
the core's eigenvalue on the one-dimensional kernel and bound
the norms of the two remaining core blocks by $M$.
Then $4\mu\sigma>M+M^2/\kappa$ gives a positive Schur
complement. Every constant here belongs to $\mathbb R$.
The exact state Gram is bounded and has this positive residue.

Finally suppose $v$ is unbounded, including the projective
endpoint. Set $w=1/v$, homogenizing before setting $w=0$.
The four event forms become
\[
 W-wX+uY-uwZ,\quad uW-Y,\quad
 uwW+uX-wY-Z,\quad Y-wZ,
\]
with bounded positive factors $D_w$ or $D_u$ where appropriate.
At $w=0$ their common kernel is
$(W,X,Y,Z)=(0,X,0,aX)$, on which the homogenized core is
$2(1+a^2)X^2$. The same real Schur argument gives a positive
residue Gram for a sufficiently large $\mu$.
For finite unbounded $v$, undo homogenization by the pure
positive testing square $v^2$; for the projective endpoint the
homogeneous expression is already the required parent.
These three cases exhaust every projective Bob test when $u$ is bounded.

It remains to consider $p=1/u$ infinitesimal, after homogenizing
by $p^2$. If $v$ is a bounded unit, use the four roots
\[
 S=pY,\qquad T=vY-Z,\qquad W,\quad X.
\]
The exact event amplitudes are
\[
 W-S,\quad J=W+vX-D_vS+pvT,\quad
 M=pvW-pX+T,\quad T.
\]
The signed core also has bounded coefficients in these roots:
\[
 2J^2+2D_pT^2+2D_v
 [p^2X^2+2vSX-2pXT-v^2S^2+2pvST-p^2T^2].
\]
At $p=0$, set $b=\bar v\ne0$, $A_0=W-S$, and
$B_0=X-bS$. The full capped residue is
\[
 (8\mu+2)T^2+4\mu(1+b^2)A_0^2
 +(4\mu+2)(A_0+bB_0)^2
 +4(1+b^2)bSB_0+2(1+b^2)b^2S^2.
\]
As $\mu$ grows, the first three positive penalties constrain
$T=A_0=B_0=0$, and the residual kernel has positive coefficient
$2(1+b^2)b^2S^2$. The quantitative real Schur argument above
therefore gives a finite $\mu$ with a strictly positive residue.

If $u,v$ are both projectively unbounded, put $w=1/v$ and
homogenize by $p^2w^2$. Use $P=Y-wZ$, $S=pZ$, $W,X$.
Now the exact event amplitudes are
\[
 M=pW-pwX+P,\quad W-pP-wS,\quad
 J=wW+X-pwP-D_wS,\quad P.
\]
The exact core is
\[
 2J^2+2D_pP^2+2D_w[p^2X^2+2SX-S^2].
\]
At $p=w=0$ the capped residue is exactly
\[
 4\mu W^2+(8\mu+2)P^2+4\mu(X-S)^2+2X^2,
\]
which is strict for every $\mu>0$. The literal projective
endpoints cause no division. The literal critical axes also have direct exact certificates:
\begin{align*}
 F|_{p=0}={}&4nD_vX^2+(4m+4\alpha+2)(vY-Z)^2
 +4CD_vW^2+4AD_vY^2+(4D+2)(W+vX)^2,\\
 F|_{v=0}={}&4nD_pX^2+(4m+2)(pX+Z)^2
 +(4C+4D+2)(W-pY)^2+4AD_pY^2+4\alpha D_pZ^2.
\end{align*}
Thus only $p,v$ both nonzero infinitesimal remain untreated.
Applying the real-state result separately to real and imaginary
state components proves the complex-state claim.
\end{proof}

For a strong C2 HI source, the exact source-entry lemma below proves
that all four coefficients in the theorem are unbounded. Hence the
theorem applies without any additional source-shape assumption.
Its precise uncovered complement is the critical testing corner
$1/u$ and $v$ both nonzero infinitesimal. Exact sign gauges may
make both positive.
The original HI source leaf is therefore narrowed but not closed.
